\documentclass[../../main.tex]{subfiles}
\begin{document}

{\bf [解答]}

複素数$|\alpha|=1$と，$z$に関する方程式
\begin{align}
  z + \alpha \bar{z} + \beta = 0 \label{eq:1}
\end{align}
について考える．

\subparagraph{必要性}

\cref{eq:1}を満たす$z$が存在すると仮定する．
\cref{eq:1}の複素共役をとって
\begin{align}
  \bar{\alpha} z + \bar{z} + \bar{\beta} = 0 \label{eq:2}
\end{align}
だから，\cref{eq:1}-$\alpha\times$\cref{eq:2}より
\begin{align}
  (1-\bar{\alpha}\alpha)z + (\beta-\alpha\bar{\beta}) = 0 \\
  \alpha \bar{\beta} = \beta \quad (\because |\alpha|=1)
\end{align}
である．よって示された．

\subparagraph{十分性}

$\beta = \alpha \bar{\beta}$ の時,
$\beta = 0$ ならば $z = 0$ が\cref{eq:1}をみたす.

そこで以下$\beta \neq 0$ で考える．
\begin{align}
  & z + \alpha \bar{z} + \beta = 0 \\
  \iff
  & z + \dfrac{\beta}{\bar{\beta}} \bar{z} + \beta = 0 \\
  \iff
  & \bar{\beta} z + \beta \bar{z} + |\beta|^2 = 0 \label{eq:3}
\end{align}
である．
\cref{eq:3}は複素平面上で直線を表すからこれを満たす$z$は必ず存在する．
よって $z + \alpha \bar{z} + \beta = 0$ をみたす $z$ は必ず存在して十分である.

\newpage

\end{document}
